\subsection{An essential secondary Kodaira component}

\begin{statement}{C.2.UN.21}{Proved}{A free half-period quotient realizes secondary Kodaira normalization}
Let $q=e^{-2\pi}$ and $E=\C^*/q^{\Z}$, and use the
origin of its group structure. There is a smooth compact
threefold $Y\to\PP^1$ whose smooth general fibre is
$E\times E$ and whose two special full fibres are
\[
                        2P+2Q\quad\hbox{and}\quad 2B.
                                                            \tag{C.SKD.1}
\]
The surface $P$ is a secondary Kodaira surface, $Q$ is
a primary Hopf surface, and they meet transversely along
a smooth elliptic curve $J$. The surface $B$ is smooth
bielliptic. Both reduced special supports have Euler
number zero. In particular $P$ occurs as an essential
component in a zero-Euler torus filling.

The canonical bundle of $P$ has exact order two. For
$S=P\cup Q$, the rational cohomology ranks and coherent
dimensions are
\[
 \begin{array}{c|c|c}
 &\dim H^k(-,\Q)&h^k(\OO)\\ \hline
 P&(1,1,0,1,1)&(1,1,0)\\
 Q&(1,1,0,1,1)&(1,1,0)\\
 J&(1,2,1)&(1,1)\\
 S&(1,1,1,3,2)&(1,1,0)\\
 B&(1,2,2,2,1)&(1,1,0)\\
 Y&(1,1,2,4,2,1,1)&(1,1,0,0).
 \end{array}                                             \tag{C.SKD.2}
\]
The full schemes in (C.SKD.1), in contrast to their
reductions, each have coherent dimensions $(1,2,1)$.
The primary Kodaira double cover of $P$ is recovered
by the actual quadratic base change at the first
special value.
\end{statement}
\begin{proof}
Start with the compact version of C.2.UN.3 at $n=1$:
\[
 V=\operatorname{Bl}_{(0,0)}(E\times\PP^1_t),\qquad
 L=\OO_V(R),\qquad
 \widetilde Y=L^\times/q^{\Z},
                                                            \tag{C.SKD.3}
\]
where $R$ is the exceptional curve. Let $\sigma$ be
the canonical section of $L$. The projection
$\widetilde Y\to V\to\PP^1_t$ is proper and smooth
off its reduced central divisor
$\widetilde P+\widetilde Q$. The first component is
the primary Kodaira surface
$\OO_E([0])^\times/q^{\Z}$, the second is the Hopf
surface $\OO_{\PP^1}(-1)^\times/q^{\Z}$, and their
intersection is an elliptic curve. These are the
actual line-bundle quotient maps of C.2.UN.3.

The involution $(p,t)\mapsto(-p,-t)$ fixes the blowup
centre and lifts to an involution $\iota$ of $V$.
It has the canonical order-two lift $\widehat\iota$
to $L$ preserving $\sigma$. Define an automorphism
of $\widetilde Y$ by
\[
             g[\ell]=[q^{1/2}\widehat\iota(\ell)].
                                                            \tag{C.SKD.4}
\]
Its square is multiplication by $q$, so $g^2=1$
on $\widetilde Y$. It acts freely. Indeed at a
fixed base point the order-two linearization on
the line $L_v$ has eigenvalue $+1$ or $-1$.
Neither $q^{1/2}$ nor $-q^{1/2}$ belongs to
$q^{\Z}$. Thus (C.SKD.4) has no fixed point on
any elliptic fibre, including those over $R$.
The quotient
\[
              \pi:\widetilde Y\longrightarrow Y
                        =\widetilde Y/\langle g\rangle
                                                            \tag{C.SKD.5}
\]
is finite \'etale of degree two, and $Y$ is smooth
and compact. The function $T=t^2$ defines its map
to $\PP^1_T$. For $T\ne0,\infty$, the two values
of $t$ are interchanged by $g$, so each quotient
fibre is exactly one of the original product
tori.

The strict central base curve and $R$ are
invariant. The local equation upstairs at their
intersection is $t=xy$, hence $T=x^2y^2$.
Since (C.SKD.5) is \'etale, this proves the first
full fibre and the normal-crossing assertion
in (C.SKD.1). At infinity the original smooth
fibre is $E\times E$, and the action is base
inversion on the first elliptic factor and
translation by the class of $q^{1/2}$ on the
second. It is free. Its quotient is a
bielliptic surface $B$, and the transverse
coordinate is squared, giving the full fibre
$2B$.

We identify the central quotients more precisely.
The canonical linearization on
$L|_R=\OO_{\PP^1}(-1)$ acts by $-1$ on its
tautological vector: the derivative at the
blown-up point is the scalar $-I$ on $\C^2$.
Thus the quotient of $\widetilde Q$ is
\[
                 Q=(\C^2\setminus0)/
                         (-q^{1/2})^{\Z},              \tag{C.SKD.6}
\]
again a primary Hopf surface. On $\widetilde P$
the involution covers inversion on $E$.
The nowhere-zero canonical form of the primary
Kodaira surface is the wedge of the base
one-form and the fibre logarithmic differential.
The first changes sign; the second does not.
Consequently $g$ acts by $-1$ on that form.
Its square descends to a nowhere-zero section
of $K_P^{\otimes2}$, while no canonical form
descends. Thus $K_P$ has exact order two.
The finite \'etale cover $\widetilde P\to P$
and the computed $b_1(P)=1$ below identify $P$
as a secondary Kodaira surface.

We give the topology through actual mapping
tori. Choose an $\iota$-invariant Hermitian
metric on $L$, by averaging. Let $M_5$ be its
unit circle bundle. The positive-real norm
coordinate gives
$\widetilde Y\simeq M_5\times S^1$. In that
description (C.SKD.4) is
$(x,u)\mapsto(\iota_Mx,u+\tfrac12)$, where
the period of $u$ is one. Hence the quotient
is the mapping torus
\[
          Y\simeq T(\iota_M)
            =(M_5\times\R)/
               ((x,u+1)\sim(\iota_Mx,u)).               \tag{C.SKD.7}
\]
For example the map $[x,u]\mapsto[x,2u]$
identifies the two quotient descriptions.

The integral circle-bundle Gysin sequence
over $V$ has Euler class $[R]$. The base has
free cohomology of ranks $(1,2,3,2,1)$.
Cup product with $[R]$ has ranks $1,0,1$
in degrees zero, one and two; the nonzero
images are primitive because $R^2=-1$.
Thus $M_5$ has free homology of ranks
\[
                         (1,2,4,4,2,1).                \tag{C.SKD.8}
\]
Its involution acts as $-I$ on $H_1$ and,
by oriented Poincar\'e duality, on $H_4$.
On $H_2$ it has an integral diagonal
decomposition with two $+1$ and two $-1$
summands. To specify that decomposition,
lift $E\times\{\infty\}$ and
$\{p_*\}\times\PP^1$, where $p_*\ne0$ is
a two-torsion point, using the normalized
canonical section of $L$. These are
invariant cycles, avoid $R$, and generate
the two base summands. The preimages of
the two elliptic-base loops at infinity
are circle tori; they generate the two
remaining summands and are sent to their
negatives. The homological Gysin sequence
identifies these four cycles with an
integral basis. Duality gives the same
two-positive, two-negative decomposition
on $H_3$. The actions on $H_0,H_5$ are
the identity. The Wang sequence of
(C.SKD.7) therefore gives the rational
$Y$ row of (C.SKD.2).

This also determines its integral groups,
not only their ranks. Writing $T_2=(\Z/2)^2$,
the integral cohomology of $Y$ is
\[
 \begin{array}{c|ccccccc}
 k&0&1&2&3&4&5&6\\ \hline
 H^k(Y,\Z)&
 \Z&\Z&\Z^2\oplus T_2&\Z^4\oplus T_2&
 \Z^2\oplus T_2&\Z\oplus T_2&\Z .
 \end{array}                                             \tag{C.SKD.9}
\]
Indeed the Wang maps are zero or multiplication
by $-2$ on the stated integral summands.
Their kernel terms are free, so the short
exact sequences split as abelian groups;
the universal coefficient sequence gives
(C.SKD.9).

On $\widetilde P$, the unit circle bundle
$M_3\to E$ has Euler number one and free
homology ranks $(1,2,2,1)$. Its involution
acts by $-I$ in degrees one and two and
by $+1$ in degrees zero and three.
The quotient $P$ is its mapping torus.
The same Wang maps give
\[
 \begin{array}{c|ccccc}
 k&0&1&2&3&4\\ \hline
 H^k(P,\Z)&
 \Z&\Z&T_2&\Z\oplus T_2&\Z .
 \end{array}                                             \tag{C.SKD.10}
\]
In particular its rational ranks are those
of (C.SKD.2). The Hopf ranks follow from
(C.SKD.6), which is differentiably
$S^3\times S^1$.

For the support use the exact union
$S=P\cup Q$ with intersection $J$. The
radial deck loop of $J$ maps to the
mapping-torus generator on $P$ and the
Hopf generator on $Q$. Its angular loop
maps to zero in rational first homology:
upstairs it is the circle of an
Euler-number-one bundle, hence is already
zero in the abelianization. Therefore
the map
\[
           H^1(P,\Q)\oplus H^1(Q,\Q)
                         \longrightarrow H^1(J,\Q)
                                                            \tag{C.SKD.11}
\]
has rank one. Both surfaces have zero
rational $H^2$, so the next restriction
map has rank zero. Mayer--Vietoris gives
the support ranks $(1,1,1,3,2)$.
There is also an exact integral comparison at the only
potential torsion restriction. In the unit circle bundle
$M_3\to E$, the core fibre over zero bounds an
$\iota_M$-invariant oriented surface. To construct it,
take the normalized canonical section over the complement
of a small invariant disc at zero. Choose a local
coordinate and an equivariant frame with
$\sigma=z e$; the lift is
$(z,u)\mapsto(-z,-u)$ in the unit-circle coordinates.
The section on the disc boundary is
$(z,u)=(\varepsilon e^{i\theta},e^{i\theta})$.
Attach the annulus
\[
 (r,\theta)\longmapsto(re^{i\theta},e^{i\theta}),
                 \qquad 0\le r\le\varepsilon .
\]
Its inner boundary is the core circle. The involution
is $\theta\mapsto\theta+\pi$ on this annulus and
preserves its orientation; it also preserves the
orientation of the section over the punctured elliptic
curve. The joined surface, or its invariant rounded
chain, therefore has an oriented mapping torus whose
boundary is $J$. Thus
$H_2(J,\Z)\to H_2(P,\Z)$ is zero integrally,
including its possible torsion target. The angular
first-homology class is already zero in $M_3$;
the other generator maps to a primitive mapping-torus
class on $P$ and to a primitive Hopf class on $Q$.
The integral Mayer--Vietoris sequence now gives
\[
 \begin{array}{c|ccccc}
 k&0&1&2&3&4\\ \hline
 H_k(S,\Z)&
 \Z&\Z\oplus T_2&\Z\oplus T_2&\Z^3&\Z^2\\
 H^k(S,\Z)&
 \Z&\Z&\Z\oplus T_2&\Z^3\oplus T_2&\Z^2 .
 \end{array}                                             \tag{C.SKD.14}
\]
The free kernel terms split as abelian groups.
Any torsion summand in the radial image is absorbed
by eliminating the primitive Hopf generator, so
no choice of radial lift changes these groups.
The invariant cohomology of the free
quotient of $E\times E$ at infinity
gives the row for $B$. In both cases
the displayed Euler number is zero.
For completeness, $B$ is also the product of an
angular circle with the mapping torus of inversion
on the first elliptic curve. The Wang maps on
that curve are $0,-2I_2,0$. Its integral
cohomology is therefore
\[
 H^k(B,\Z)=
 \bigl(\Z,\ \Z^2,\ \Z^2\oplus T_2,\
                  \Z^2\oplus T_2,\ \Z\bigr)_k .
                                                            \tag{C.SKD.15}
\]
The integral groups of $J$ and $Q$ are free with
the ranks already displayed.

For coherent cohomology, the elliptic
line-bundle quotient $\widetilde Y\to V$
has direct images $\OO_V$ in degrees
zero and one and zero above one. The
degree-one frame is the constant
logarithmic deck cocycle, as in
(C.PE.21). Its involution action is
$+1$, while the action on
$H^1(V,\OO_V)=H^1(E,\OO_E)$ is $-1$.
The finite \'etale invariant-sheaf
comparison for (C.SKD.5), applied to
Leray's filtration, gives
$h^k(Y,\OO_Y)=(1,1,0,0)$. The identical
calculation for $\widetilde P\to E$
gives $h^k(P,\OO_P)=(1,1,0)$, consistent
with the canonical character.
The Hopf and elliptic rows were
computed from their actual deck cocycles
in C.2.UN.18. In particular
$H^1(Q,\OO_Q)\to H^1(J,\OO_J)$ is
nonzero and hence onto. The coherent
normalization sequence gives
$h^k(S,\OO_S)=(1,1,0)$. The quotient
at infinity has the same base-negative,
fibre-positive action, giving
$h^k(B,\OO_B)=(1,1,0)$.

We finally verify the full scheme fibres.
Upstairs let $\widetilde S$ be the
reduced central fibre of the original
map to $t$. Its coherent normalization
sequence gives $(1,2,1)$, the same
dimensions as a nearby product torus.
Grauert's theorem therefore makes all
$R^qf_*\OO_{\widetilde Y}$ locally free
near $t=0$, of ranks $1,2,1$, with
base change. The actual inverse image
of the full fibre $2S$ under the
\'etale map (C.SKD.5) is $2\widetilde S$,
defined by $t^2=0$. Thus
\[
 H^q(2S,\OO_{2S})
  \simeq
 \left((R^qf_*\OO_{\widetilde Y})/
                   t^2(R^qf_*\OO_{\widetilde Y})\right)^g .
                                                            \tag{C.SKD.12}
\]
Choose a basis of eigenvectors on the
central fibre of each direct-image
bundle. Lift each vector and average
with $g$ with its specified sign; the
resulting holomorphic sections remain
a frame near zero and satisfy
$g(e)=\pm e$. Since $g(t)=-t$, the
invariants in each summand
$\C[t]/(t^2)e$ have dimension one:
they are generated by $e$ for sign
$+1$ and by $te$ for sign $-1$.
Equation (C.SKD.12) gives precisely
$(1,2,1)$. The same argument at
infinity proves the assertion for
$2B$. This comparison retains both
the nilpotent direction and its sign.

No rational curve lies in $P$:
its map from $\PP^1$ would lift to
the primary Kodaira cover, whose
elliptic base and elliptic fibres
force that map to be constant.
The same lifting argument to
$\C^2\setminus0$ treats $Q$.
The inverse-map resolution comparison
of C.2.UN.3 therefore proves their
essentiality. At $T=0$, the finite
map
\[
 \widetilde Y\longrightarrow
             Y\times_{\PP^1_T}\PP^1_t,\qquad
 y\longmapsto(\pi(y),f(y)),\quad T=t^2,
                                                            \tag{C.SKD.13}
\]
is generically one-to-one and has
normal source. It identifies
$\widetilde Y$ with the normalization
of the reduced fibre product.
Consequently the primary Kodaira
component is recovered by this
actual quadratic base change.
The construction supplies an
explicit secondary Kodaira component;
no birational persistence of its
secondary type under that base change
is asserted.
\end{proof}
